\section{识别与结论边界}\label{sec:identification}
本文机制使用Stackelberg\emph{时序}而不声称闭式均衡：设计者承诺$m$，追随者再响应。AEIC是报告接口，而经验对象是黑盒策略诱导的行动分布$\pi(a\mid x,m,h)$。

\paragraph{追随者响应。}在配对任务上改变$m$并固定问题、Seed、面板、轮数和信息策略，可识别公开规则对已实现模型策略的意向处理效应。正的高资源行动差异证明规则响应，但不证明任意Agent效用最大化或激励相容。

\paragraph{仅结果控制。}Static AEIC与Outcome Only区分过程可观测性和一般奖励语言。它支持局部可检查接口是有效组成，但不隔离四个AEIC字段的独立贡献；信号消融只是描述性边界。

\paragraph{剂量响应。}式~\eqref{eq:followerutility}预测奖励提高时高资源行动弱增加、成本提高时弱降低。本文检验方向含义而不拟合结构效用模型；冻结高低点和开发中心点支持单调性，而非第二个确认性效应。

\paragraph{外部有效性。}Prompt介导的成本和信用具有后端语义，一个模型上的合同未必迁移到另一模型。因此本文将留出后端作为范围检验，不主张一般均衡、Leader全局最优、货币激励相容或拓扑无关有效性。

\paragraph{可证伪预测。}机制给出三个先验预测：其一，若Agent读取并响应过程合同，提高过程信用应使$\Pr(H(a)=1)$不下降；其二，提高同一行动的显式成本应使该概率不升高；其三，若效应来自可局部满足的生产规则而非一般奖励措辞，Outcome Only不应复制Static AEIC的行动响应。实验分别用奖励扫描、成本扫描和三条件确认性比较检验这些预测。预测失败将直接削弱机制解释，而不是被事后改写为提示效果。

\paragraph{可识别与不可识别部分。}确认性处理同时包含公开信用文本、行动选择Prompt和按行动分级的执行路径，因而能够识别这套可执行合同的总体效果，不能把准确率变化归因于单个字段。等计算量对照检验“只是多给Token”的替代解释；报告接口实验检验AEIC相对自由文本及一般结构化输出的增量；两者共同缩小解释范围，但仍不构成四字段的独立因果分解。
